\chapter{Experiments}\label{chapter-4-experiments}

\section{Experimental Design}\label{41-experimental-design}

\subsection{Research Questions}\label{411-research-questions}

This experimental study addresses the following research questions:

\begin{longtable}[]{@{}lll@{}}
\toprule\noalign{}
RQ & Question & Metric(s) \\
\midrule\noalign{}
\endhead
\bottomrule\noalign{}
\endlastfoot
RQ1 & How do travel-time distributions compare across simulators for
identical inputs? & RMSE, GEH, KS \\
RQ2 & How does runtime scale with network size and demand volume across
engines? & Wall-clock, throughput \\
RQ3 & How consistent are results across repeated runs with different
random seeds? & \(R = 1 - \sigma/\mu\) \\
RQ4 & What are the practical trade-offs between fidelity, speed, and
reproducibility? & Multi-metric analysis \\
\end{longtable}

\subsection{Independent
Variables}\label{412-independent-variables}

\textbf{Variable 1: Simulator Engine (4 configurations evaluated)}

\begin{longtable}[]{@{}llll@{}}
\toprule\noalign{}
Engine & Paradigm & Traffic Model & Hardware \\
\midrule\noalign{}
\endhead
\bottomrule\noalign{}
\endlastfoot
SUMO (micro) & Microscopic & Krauss car-following + lane-changing &
CPU \\
SUMO (meso) & Mesoscopic & Queue-based link traversal & CPU \\
MATSim & Mesoscopic & Activity-based, event-driven queues & CPU \\
DTALite & Mesoscopic & Dynamic Traffic Assignment (UE) & CPU \\
\end{longtable}

\begin{quote}
The plan listed five engines (SUMO, MATSim, POLARIS, LPSim,
QarSUMO). After two integration cycles, the matrix narrows to
\textbf{three primary engines} chosen for paradigm spread: SUMO
microscopic (car-following), MATSim queue-based agent simulation, and
DTALite mesoscopic Dynamic Traffic Assignment. Three of the
originally-proposed engines were systematically evaluated and ruled out:
QarSUMO dropped in Version\_4 Phase A (no usable public source), LPSim
integrated in Version\_4 Phase B and abandoned in Version\_5 (GPU kernel
SIGSEGV on networks \textgreater{} a few-K nodes), POLARIS and CityFlow
evaluated as third-engine alternatives and rejected at criteria. Full
retrospectives in \href{../engines/}{\texttt{doc/engines/}}. The
canonical 11-cell matrix exercises all four (engine × mode) combinations
across three scenarios at N = 5 repeats each, and is fully CPU-only ---
the 1 K tier runs end-to-end on a Mac laptop in under a minute; the 10 K
and 50 K tiers fit inside a single Pitzer SLURM job
(\textasciitilde22-23 h wall-clock) with the Phase 12 BFS-prep cache.
\end{quote}

\textbf{Variable 2: Scenario Scale}

The canonical \texttt{benchmark\_small.yaml} runspec covers
\textbf{three tiers} spanning two orders of magnitude in trip count. The
1 K tier is laptop-runnable; the 10 K and 50 K tiers require an HPC node
for the SUMO micro and DTALite cells. Larger tiers (200 K, 500 K) are
generated locally via helper scripts and reserved for separate HPC runs.

\begin{longtable}[]{@{}lllll@{}}
\toprule\noalign{}
Tier & Trip Count & Helper script & In \texttt{benchmark\_small.yaml} &
Where it runs \\
\midrule\noalign{}
\endhead
\bottomrule\noalign{}
\endlastfoot
\textbf{Bundled} & \textbf{1,000} &
\texttt{scripts/01\_chicago\_1k\_car.py} & ✅ & Laptop (Apple
Silicon) \\
\textbf{Small} & \textbf{10,000} & \texttt{scripts/02\_nyc\_10k\_car.py}
& ✅ & Laptop (cached) / cluster \\
\textbf{Medium} & \textbf{50,000} & \texttt{scripts/03\_la\_50k\_car.py}
& ✅ & Cluster (Pitzer 36 h walltime) \\
Large & 200,000 & \texttt{scripts/04\_chicago\_200k\_car.py} & (separate
runspec) & Cluster \\
Stress & 500,000 & \texttt{scripts/05\_nyc\_500k\_car.py} & (separate
runspec) & Cluster (Linux only) \\
\end{longtable}

\textbf{Variable 3: City Network (3 metropolitan areas in the canonical
matrix)}

\begin{longtable}[]{@{}llllll@{}}
\toprule\noalign{}
City & Bundle & Nodes (in SCC) & Links (in SCC) & SCC \% & Character \\
\midrule\noalign{}
\endhead
\bottomrule\noalign{}
\endlastfoot
Chicago & \texttt{chicago\_1k\_car} & 19,744 & 58,432 & 98.4 \% & Dense
grid, mixed use \\
NYC & \texttt{nyc\_10k\_car} & 33,988 & 105,009 & 99.6 \% &
Multi-borough mix \\
LA & \texttt{la\_50k\_car} & 158,690 & 466,518 & 99.8 \% & Sprawling
freeway-dominated \\
\end{longtable}

The three scenarios span \textasciitilde8× in node count and
\textasciitilde8× in link count, while the trip count spans 50× (1K →
50K). This decouples the \emph{network-size} effect from the
\emph{demand-density} effect in the scaling analysis (§5.1, §5.5).

\textbf{Why Chicago, NYC, and LA:}

\begin{itemize}
\tightlist
\item
  \textbf{Chicago} --- classic American grid; mix of arterials and
  residential streets; ModelGen census coverage is most complete. Dense
  enough to exercise all three engines end-to-end while remaining small
  enough to ship as a bundled artefact.
\item
  \textbf{NYC} --- multi-borough, irregular network topology; tests the
  engines on a network with mixed grid + organic geometry. The 10 K trip
  volume crosses the threshold where SUMO micro becomes a coffee-break
  run (\textasciitilde3 min per seed).
\item
  \textbf{LA} --- sprawling freeway-dominated network; largest tested
  network (\textasciitilde470 K links). The 50 K trip volume tests the
  engines at the boundary of practical UE convergence (DTALite hits a
  path4gmns binary-side scaling ceiling at this size --- see §5.7).
\end{itemize}

\subsection{Dependent Variables}\label{413-dependent-variables}

\textbf{Fidelity (RQ1):}

\begin{longtable}[]{@{}llll@{}}
\toprule\noalign{}
Metric & Formula & Unit & Interpretation \\
\midrule\noalign{}
\endhead
\bottomrule\noalign{}
\endlastfoot
Mean travel time & \(\bar{t} = \frac{1}{n}\sum_i t_i\) & seconds &
Central tendency of trip durations \\
P95 travel time & 95th percentile of \(\{t_i\}\) & seconds & Tail
behaviour / worst-case trips \\
Trip completion rate & \(n_\text{completed} / n_\text{feasible}\) & \% &
How many feasible trips reached destination \\
RMSE (cross-simulator) & \(\sqrt{\frac{1}{n}\sum(y_i - \hat{y}_i)^2}\) &
seconds & Agreement between two simulators \\
GEH (link-level) & \(\sqrt{\frac{2(M-C)^2}{M+C}}\) & dimensionless &
Traffic engineering fit metric \\
KS statistic & \(\sup_x |F_1(x) - F_2(x)|\) & dimensionless &
Distribution shape comparison \\
\end{longtable}

\textbf{Scalability (RQ2):}

\begin{longtable}[]{@{}llll@{}}
\toprule\noalign{}
Metric & Formula & Unit & Interpretation \\
\midrule\noalign{}
\endhead
\bottomrule\noalign{}
\endlastfoot
Wall-clock runtime & \texttt{perf\_counter()} end -- start & seconds &
Absolute performance \\
Throughput & trips\_completed / runtime & trips/sec & Efficiency
metric \\
SRT & simulated\_time / wall\_clock\_time & ratio &
Faster-than-real-time indicator \\
Peak memory & OS-reported peak RSS & MB & Hardware requirement \\
\end{longtable}

\textbf{Reproducibility (RQ3):}

\begin{longtable}[]{@{}llll@{}}
\toprule\noalign{}
Metric & Formula & Range & Interpretation \\
\midrule\noalign{}
\endhead
\bottomrule\noalign{}
\endlastfoot
Reproducibility index \(R\) & \(1 - \sigma/\mu\) & {[}-∞, 1.0{]} & 1.0 =
perfectly deterministic \\
Max deviation & \(\max_i |x_i - \mu|\) & same as KPI & Worst-case
variability \\
CV (coefficient of variation) & \(\sigma/\mu\) & {[}0, ∞) & Relative
variability (lower = better) \\
\end{longtable}

\subsection{Canonical Experimental
Matrix}\label{414-canonical-experimental-matrix}

The matrix declared by \texttt{runspecs/benchmark\_small.yaml} (11 cells
× 5 repeats = \textbf{55 runs total}):

\begin{longtable}[]{@{}lllll@{}}
\toprule\noalign{}
Scenario & Engine & Mode & Trips & Repeats \\
\midrule\noalign{}
\endhead
\bottomrule\noalign{}
\endlastfoot
chicago\_1k\_car & sumo & mesoscopic & 1,000 & 5 \\
chicago\_1k\_car & sumo & microscopic & 1,000 & 5 \\
chicago\_1k\_car & matsim & mesoscopic & 1,000 & 5 \\
chicago\_1k\_car & dtalite & mesoscopic & 1,000 & 5 \\
nyc\_10k\_car & sumo & mesoscopic & 10,000 & 5 \\
nyc\_10k\_car & sumo & microscopic & 10,000 & 5 \\
nyc\_10k\_car & matsim & mesoscopic & 10,000 & 5 \\
nyc\_10k\_car & dtalite & mesoscopic & 10,000 & 5 \\
la\_50k\_car & sumo & mesoscopic & 50,000 & 5 \\
la\_50k\_car & matsim & mesoscopic & 50,000 & 5 \\
la\_50k\_car & dtalite & mesoscopic & 50,000 & 5 \\
\end{longtable}

la\_50k\_car drops SUMO microscopic by design --- at 50 K trips a single
arm64 SUMO micro run wall-clocks past 4 h, which is impractical for a
"small tier" matrix and provides little additional paradigm signal
beyond chicago\_1k + nyc\_10k micro. MATSim and DTALite are
mesoscopic-only by design.

Every \texttt{feasibility\_report.json} confirms
\texttt{feasible\_trips} matches the trip count for every adapter
(verified by Q1 of \texttt{audit\_fairness}), so the four engines
simulate exactly the same trip set per cell. All three engines are
CPU-only and the full matrix completes inside a single Pitzer SLURM job
(\textasciitilde22-23 h with the Phase 12 BFS-prep cache populated by
the first seed of each scenario × engine pair).

\textbf{Control variables} (held constant across all conditions):

\begin{longtable}[]{@{}lll@{}}
\toprule\noalign{}
Parameter & Value & Rationale \\
\midrule\noalign{}
\endhead
\bottomrule\noalign{}
\endlastfoot
Random seeds & 42, 43, 44, 45, 46 & Reproducible; 5 runs per cell \\
Demand strategy & Census-calibrated (ModelGen + PUMS) & Consistent
realistic inputs \\
Routing & State-aware BFS with OSM turn restrictions & Deterministic,
version-independent \\
Time horizon & 1 hour (3,600 s) & Standard morning peak period \\
MATSim iterations & 1 (\texttt{lastIteration\ =\ 0}, no replanning) &
Fair comparison with SUMO single-pass mode \\
DTALite iterations & 5 column-gen + 5 column-update & Standard UE
convergence \\
Feasibility filter & SCC-based & Same trip set fed to every engine \\
\end{longtable}

\begin{center}\rule{0.5\linewidth}{0.5pt}\end{center}

\section{Scenario Descriptions}\label{42-scenario-descriptions}

\subsection{\texorpdfstring{4.2.1
\texttt{chicago\_1k\_car}}{4.2.1 chicago\_1k\_car}}\label{421-chicago_1k_car}

\textbf{Generated}: \texttt{python\ scripts/01\_chicago\_1k\_car.py}
(equivalent to
\texttt{python\ generate.py\ -\/-city\ chicago\ -\/-trips\ 1000\ -\/-modes\ car\ -\/-seed\ 42})

\begin{longtable}[]{@{}ll@{}}
\toprule\noalign{}
Property & Value \\
\midrule\noalign{}
\endhead
\bottomrule\noalign{}
\endlastfoot
Nodes & 1,245 \\
Links & 2,862 \\
Largest SCC & 1,204 nodes (96.7 \%), 2,796 links \\
Bounding box & 41.842°--41.914° N, -87.678°---87.582° W \\
Radius from centre & 4 km \\
Source & OpenStreetMap (cached in \texttt{cache/}) \\
Total trips & 1,000 \\
Demand strategy & Census-calibrated (ModelGen) \\
Mode & 100 \% car \\
Time window & 0 -- 3,600 s (one-hour morning peak) \\
Census commuters in bbox & \textasciitilde15,000 -- 25,000 \\
\end{longtable}

\textbf{Key features}: Dense grid pattern with regular block spacing;
high intersection density; mixed residential/commercial; multiple
highway segments (I-90/94, I-290).

\subsection{Larger Tiers (HPC)}\label{422-larger-tiers-hpc}

The 10K -- 500K tiers are generated by
\texttt{scripts/02\_nyc\_10k\_car.py},
\texttt{scripts/03\_la\_50k\_car.py},
\texttt{scripts/04\_chicago\_200k\_car.py}, and
\texttt{scripts/05\_nyc\_500k\_car.py}. They target NYC, LA, and Chicago
respectively and run on OSC Pitzer --- see
\href{../PITZER.md}{\texttt{doc/PITZER.md}} for the operational guide.

\begin{center}\rule{0.5\linewidth}{0.5pt}\end{center}

\section{Execution Environment}\label{43-execution-environment}

\subsection{Development Environment
(Local)}\label{431-development-environment-local}

\begin{longtable}[]{@{}ll@{}}
\toprule\noalign{}
Component & Specification \\
\midrule\noalign{}
\endhead
\bottomrule\noalign{}
\endlastfoot
Machine & MacBook Pro (Apple Silicon M4 Pro) \\
CPU & 14 cores (10P + 4E) \\
RAM & 16 GB \\
Storage & 1 TB SSD \\
GPU & Integrated (no discrete NVIDIA) \\
OS & macOS 14.x \\
\end{longtable}

\subsection{HPC Environment (OSC
Pitzer)}\label{432-hpc-environment-osc-pitzer}

Used for the canonical 11-cell \texttt{benchmark\_small} matrix
(chicago\_1k + nyc\_10k + la\_50k). The 1 K tier alone is
laptop-runnable; the 10 K and 50 K tiers need HPC because (a) SUMO micro
at 10 K is a coffee-break run, and (b) DTALite at 10 K + has
super-linear UE iteration cost. All three primary engines (SUMO, MATSim,
DTALite) are CPU-only after the LPSim removal in Version\_5. See
\texttt{doc/PITZER.md} for the operational guide and
\texttt{cluster/jobs/benchmark\_small.sbatch} for the canonical sbatch
script (Phase 12.4: \texttt{-\/-mem=128G},
\texttt{-\/-cpus-per-task=16}, \texttt{-\/-time=36:00:00}).

\begin{longtable}[]{@{}ll@{}}
\toprule\noalign{}
Component & Specification \\
\midrule\noalign{}
\endhead
\bottomrule\noalign{}
\endlastfoot
Cluster & Ohio Supercomputer Center --- Pitzer (RHEL 9, SLURM) \\
CPU per node & Intel Xeon Skylake (40 cores) or Cascade Lake (48 cores),
192 GB \\
Large-mem nodes & Up to 3 TB RAM (\texttt{hugemem} partition) \\
Partitions used & \texttt{cpu} (all three shipped engines are CPU-only
after the Version\_5 LPSim retirement) \\
Max wall time & 7 days (\texttt{cpu}, \texttt{gpu}); 14 days
(\texttt{longcpu}, restricted) \\
Storage & Home 500 GB, Project (\texttt{PMIU0110}) 500 GB, scratch
per-job \\
Project account & \texttt{-\/-account=PMIU0110} \\
Internet access & Outbound via NAT on login + compute nodes (PyPI,
GitHub; Overpass used only as fallback --- primary OSM ingest uses state
PBFs staged under \texttt{\textasciitilde{}/SimForge/osm\_data/}) \\
\end{longtable}

\subsection{Software Versions}\label{433-software-versions}

\begin{longtable}[]{@{}llll@{}}
\toprule\noalign{}
Software & Version & Installation & Notes \\
\midrule\noalign{}
\endhead
\bottomrule\noalign{}
\endlastfoot
Python & 3.13.x & \texttt{setup\_simforge.py} & Bootstrapper \\
SUMO & 1.26.0 & \texttt{uv\ pip\ install\ eclipse-sumo==1.26.0}
(separate from \texttt{requirements.lock}; the wheel is
manylinux\_2\_28\_x86\_64-only and pinned outside the cross-platform
lockfile) & Mandatory \\
MATSim & 15.0 & JAR & \texttt{lib/matsim-15.0/} \\
Java & 17 & Homebrew & MATSim runtime \\
DTALite & path4gmns 0.10.0 & bundled in \texttt{requirements.lock}
(\texttt{path4gmns} wheel) & CPU mesoscopic DTA; on Mac needs
\texttt{brew\ install\ libomp} \\
osmnx & 2.x (\texttt{\textgreater{}=2.0,\textless{}3}) & pip & Parse PBF
slice → graph, bbox truncate (\texttt{bbox=(W,S,E,N)} tuple) \\
pyosmium & 4.x & pip (\texttt{osmium}) & Bbox-slice Geofabrik PBFs in
\texttt{osm\_data/} \\
lxml & 5.x & pip & XML processing \\
pandas & 2.x & pip & Demand CSV handling \\
\end{longtable}

\subsection{Execution Protocol}\label{434-execution-protocol}

For each cell of the 11-cell matrix:

\begin{enumerate}
\def\labelenumi{\arabic{enumi}.}
\tightlist
\item
  \textbf{Validation}: \texttt{validate\_bundle.py} verifies hash +
  referential integrity for the scenario bundle.
\item
  \textbf{Conversion}: Adapter writes simulator-specific inputs (network
  XML, plans/routes, settings) plus a \texttt{feasibility\_report.json}
  sidecar. Phase 12 BFS-prep cache: the first seed populates a
  per-engine cache; the remaining 4 seeds hardlink the cached prep into
  their cell directories (\textasciitilde93 s/seed for cached SUMO meso
  vs \textasciitilde28,681 s/seed for the cold first seed at la\_50k).
\item
  \textbf{Measurement runs}: 5 runs per cell with seeds \{42, 43, 44,
  45, 46\}.
\item
  \textbf{Metric extraction}: Parse simulator outputs → travel-time
  stats, throughput. SUMO via \texttt{parse\_sumo\_tripinfo}
  (tripinfo.xml), MATSim via \texttt{parse\_matsim\_output}
  (output\_trips.csv.gz), DTALite via \texttt{parse\_dtalite\_output}
  (link\_performance.csv + agent.csv).
\item
  \textbf{Result serialisation}:
  \texttt{runs/benchmark\_small/\textless{}scenario\textgreater{}/benchmark\_results\_benchmark\_small.json}
  per scenario; the parallel-by-scenario sbatch (Phase 12) produces
  three independent JSONs that the analyzer + audit\_fairness can read
  together.
\end{enumerate}

\begin{center}\rule{0.5\linewidth}{0.5pt}\end{center}

\section{Measured Results}\label{44-measured-results}

\subsection{Headline numbers (OSC Pitzer Intel Xeon Skylake, 16
cores per
worker)}\label{441-headline-numbers-osc-pitzer-intel-xeon-skylake-16-cores-per-worker}

The numbers below come from Pitzer SLURM jobs \texttt{47237978} (initial
run) + \texttt{47248311} (post-Phase-12.4 re-queue) +
\texttt{tools/recover\_partial\_summary.py} (Phase 12.5 synthesis for
la\_50k\_car DTALite cells). See \href{../../CHANGELOG.md}{CHANGELOG.md}
Phase 12 series for the full diagnostic chain. Chapter 5 (Results)
reproduces these numbers verbatim from the on-disk JSONs at
\texttt{runs/benchmark\_small/\textless{}scenario\textgreater{}/benchmark\_results\_benchmark\_small.json}
--- see §5.1 (Table 5.1), §5.2 (Table 5.2), §5.3 (Fig 5.3 cross-engine
ratios), §5.7 (DTALite scaling-ceiling discussion).

Compact summary (full 11-cell table is Table 5.1 in Chapter 5):

\begin{longtable}[]{@{}lllll@{}}
\toprule\noalign{}
Scenario & SUMO/MATSim mean-TT ratio & SUMO meso runtime & MATSim meso
runtime & DTALite meso runtime \\
\midrule\noalign{}
\endhead
\bottomrule\noalign{}
\endlastfoot
chicago\_1k\_car & 0.869 (-13.1 \%) & 9.50 s & 8.00 s & 22.26 s \\
nyc\_10k\_car & 1.132 (+13.2 \%) & 19.90 s & 15.24 s & 583.06 s \\
\textbf{la\_50k\_car} & \textbf{1.046 (+4.6 \%)} & 91.68 s & 52.98 s &
\emph{did not converge --- see §5.7} \\
\end{longtable}

\subsection{Fidelity (RQ1)}\label{442-fidelity-rq1}

\begin{itemize}
\tightlist
\item
  \textbf{Cross-engine alignment improves with scale.} SUMO/MATSim
  mean-TT gap narrows from ±13 \% at 1 K and 10 K → +4.6 \% at 50 K. The
  convergence is law-of-large-numbers: with more trips, per-trip
  differences between SUMO\textquotesingle s stricter insertion logic
  and MATSim\textquotesingle s earlier mobsim release average out. This
  is the central paradigm-spread finding (§5.3).
\item
  \textbf{DTALite UE underestimates travel time vs queue-based mobsim by
  \textasciitilde40-45 \%} at the scales where it converges
  (chicago\_1k, nyc\_10k). Expected behaviour: equilibrium assignment
  ignores transient congestion build-up and dissipation that
  event-driven mobsim captures.
\item
  \textbf{SUMO meso vs SUMO micro disagree by 28-73 \% on mean travel
  time} (gap grows with scale) --- micro captures intersection delays
  and queue spillback.
\end{itemize}

\subsection{Scalability (RQ2)}\label{443-scalability-rq2}

\begin{itemize}
\tightlist
\item
  \textbf{SUMO meso : MATSim meso runtime ratio} narrows with scale:
  1.19× at 1 K → 1.31× at 10 K → 1.73× at 50 K (MATSim faster at 50 K;
  JVM startup amortises).
\item
  \textbf{SUMO meso vs SUMO micro speedup} widens with scale: 1.65× at 1
  K → 9.01× at 10 K (micro becomes a coffee-break run at 10 K).
\item
  \textbf{DTALite super-linear cost.} From 22 s (1 K) → 583 s (10 K) is
  26× cost for 10× trips. Beyond 10 K the path4gmns 0.10.0
  binary\textquotesingle s 4-thread cap makes it prohibitive ---
  la\_50k\_car would need \textasciitilde25 h per seed at observed rates
  (§5.7).
\item
  \textbf{MATSim\textquotesingle s per-trip cost amortises the JVM tax
  rapidly.} \textasciitilde7 s startup + per-trip work; at 50 K, MATSim
  achieves 944 trips/s --- the highest throughput at any tested tier
  (§5.5).
\end{itemize}

\subsection{Reproducibility (RQ3)}\label{444-reproducibility-rq3}

\begin{itemize}
\tightlist
\item
  \textbf{MATSim is byte-deterministic at all scales} (R = 1.0000 across
  all 3 scenarios) with \texttt{lastIteration\ =\ 0}.
\item
  \textbf{DTALite is byte-deterministic where it converges} (R = 1.0000
  at chicago\_1k and nyc\_10k; la\_50k did not converge).
\item
  \textbf{SUMO R drops with scale but stays "Good"}: 0.9982 (1 K) →
  0.9564 (10 K, worst case) → 0.9804 (50 K). CV stays \textless{} 5 \%
  everywhere. SUMO is \emph{reproducible} in the engineering sense but
  not byte-deterministic like MATSim and DTALite.
\end{itemize}

\subsection{Trade-off Analysis
(RQ4)}\label{445-trade-off-analysis-rq4}

\begin{verbatim}
Fidelity ▲
         │  ● SUMO-micro     (highest fidelity; 9× slower than meso at 10K, infeasible at 50K)
         │
         │      ● SUMO-meso  (lower fidelity, scales linearly with trip count)
         │      ● MATSim     (different mobsim, byte-deterministic; JVM-tax floor)
         │      ● DTALite    (UE equilibrium; -40% mean TT; super-linear cost; 4-thread cap at 50K)
         │
         └──────────────────────────────────────▶ Speed
\end{verbatim}

\textbf{Key trade-off}: The 1 K tier is too small to be practically
interesting --- all four cells finish in ≤ 22 s. The decisive trade-offs
surface at 10 K + (SUMO micro becomes coffee-break) and 50 K + (SUMO
micro infeasible, DTALite hits path4gmns scaling ceiling). The published
thesis numbers cover all three regimes; future work would extend to 200
K and 500 K once the path4gmns ceiling is mitigated (§5.7).

\begin{center}\rule{0.5\linewidth}{0.5pt}\end{center}

\section{Analysis Methods}\label{45-analysis-methods}

\subsection{Statistical Analysis}\label{451-statistical-analysis}

\textbf{Cross-simulator comparison (RQ1)}

\begin{itemize}
\tightlist
\item
  Paired comparisons: each pair of simulators compared on identical
  scenarios.
\item
  KS test at α = 0.05: determines whether travel-time distributions
  differ significantly.
\item
  GEH acceptance criterion: ≥ 85 \% of links with GEH \textless{} 5.0
  indicates acceptable fit.
\end{itemize}

\textbf{Scaling analysis (RQ2)}

\begin{itemize}
\tightlist
\item
  Log-log plot of runtime vs trip count → slope indicates scaling
  exponent.
\item
  Expected: meso slopes ≈ 1.0 (linear), micro slopes ≈ 1.2 -- 1.5
  (super-linear).
\end{itemize}

\textbf{Reproducibility analysis (RQ3)}

\begin{itemize}
\tightlist
\item
  Compute \(R\) for each \texttt{(scenario,\ engine,\ mode,\ KPI)}
  combination.
\item
  Report \(R \geq 0.99\) as "effectively deterministic"; flag any
  \(R < 0.90\) for investigation.
\end{itemize}

\subsection{Visualisation Plan}\label{452-visualisation-plan}

\texttt{evaluation/generate\_plots.py} emits 10 figures (5.1 -- 5.10).
See \href{../RESULTS_GUIDE.md}{doc/RESULTS\_GUIDE.md} for the per-figure
description.

\subsection{Output Formats}\label{453-output-formats}

\texttt{evaluation/analyze\_benchmark.py} emits:

\begin{itemize}
\tightlist
\item
  \textbf{Console summary} with categorised metrics
\item
  \textbf{LaTeX tables} ready for thesis inclusion (\texttt{-\/-latex})
\item
  \textbf{Markdown tables} for documentation (\texttt{-\/-markdown})
\item
  \textbf{Coverage diagnostic} flagging low-sample / asymmetric /
  silently-failed cells
\end{itemize}

\begin{Shaded}
\begin{Highlighting}[]
\ExtensionTok{python} \AttributeTok{{-}m}\NormalTok{ evaluation.analyze\_benchmark runs/benchmark\_small/benchmark\_results\_benchmark\_small.json }\AttributeTok{{-}{-}latex} \AttributeTok{{-}{-}markdown}
\end{Highlighting}
\end{Shaded}

\subsection{Cross-engine fairness
audit}\label{454-cross-engine-fairness-audit}

\texttt{evaluation/audit\_fairness.py} is the methodology check that
verifies the cross-engine comparison was actually fair before any of its
results are reported. For each scenario, it compares the inputs each
engine consumed and the outputs each engine produced, organised as four
checks:

\begin{longtable}[]{@{}lll@{}}
\toprule\noalign{}
Check & What it verifies & Pass condition \\
\midrule\noalign{}
\endhead
\bottomrule\noalign{}
\endlastfoot
\textbf{Q1} & Cross-engine feasibility verdict & All
engines\textquotesingle{} \texttt{feasibility\_report.json}
byte-identical (same \texttt{feasible\_trips}, \texttt{total\_trips},
\texttt{scc\_nodes}, \texttt{scc\_links}, skip counts) \\
\textbf{Q2} & Engine input network & All engines emit identical
SCC-filtered node + link counts (modulo documented OSM-noise drops) \\
\textbf{Q3} & Trips actually simulated & Each engine simulates exactly
the cross-engine \texttt{feasible\_trips} count (MATSim plans, DTALite
demand-volume sum, SUMO routes) \\
\textbf{Q4} & Cross-engine travel-time spread & Per-engine mean / P95 /
completion + pairwise mean-TT ratios --- this is the paradigm-spread
signal \\
\end{longtable}

\begin{Shaded}
\begin{Highlighting}[]
\ExtensionTok{python} \AttributeTok{{-}m}\NormalTok{ evaluation.audit\_fairness runs/benchmark\_small}
\end{Highlighting}
\end{Shaded}

The audit is read-only and emits a defender-friendly report.
PASS/WARN/FAIL verdicts on Q1--Q3 are the methodology-section evidence
that "the engines were given the same problem"; the Q4 ratios are the
paradigm-spread signal at the heart of the cross-engine comparison.
Sample audit output on the bundled chicago\_1k\_car (Pitzer, all three
engines) is recorded in \texttt{doc/EXPERIMENT\_LOG.md} §3.

\begin{center}\rule{0.5\linewidth}{0.5pt}\end{center}

\section{Threats to Validity}\label{46-threats-to-validity}

\subsection{Internal Validity}\label{461-internal-validity}

\begin{longtable}[]{@{}ll@{}}
\toprule\noalign{}
Threat & Mitigation \\
\midrule\noalign{}
\endhead
\bottomrule\noalign{}
\endlastfoot
Random seed affecting results & 5 runs per condition with seeds \{42,
43, 44, 45, 46\}; report mean ± 95 \% CI \\
JVM warm-up affecting MATSim & All runs include the same JVM start cost;
comparison is fair-relative; cost amortises \textless{} 30 \% at 10 K
+ \\
OS scheduling noise & Use \texttt{perf\_counter()}; HPC runs on
dedicated nodes; cached cell std \textless{} 5 \% CV \\
Adapter conversion errors & \textasciitilde639 unit tests (with all 5
bundles generated) including byte-identical determinism tests + the
three-function adapter-contract regression guard \\
Scenario validation failures & Pre-flight validation check before every
run \\
Trip-count asymmetry across engines & SCC filter at generator + adapter;
\texttt{feasibility\_report.json} audit trail; \texttt{audit\_fairness}
Q1 PASS on all 3 scenarios \\
MATSim adapter route-format ambiguity & Phase 12.1 fix:
\texttt{\textless{}route\ type="links"\textgreater{}} text content
includes start\_link + end\_link tokens; verified by 0-trip → 1000-trip
empirical check \\
path4gmns 0.10.0 4-thread cap on DTALite & Discovered Phase 12.5;
documented as scaling ceiling (§5.7); affects la\_50k\_car DTALite cells
only \\
\end{longtable}

\subsection{External Validity}\label{462-external-validity}

\begin{longtable}[]{@{}ll@{}}
\toprule\noalign{}
Threat & Mitigation \\
\midrule\noalign{}
\endhead
\bottomrule\noalign{}
\endlastfoot
Only 1 US city in the stress test & Larger HPC tiers cover NYC at 10K
and 500K, LA at 50K, Chicago at 200K \\
Census-calibrated (not real OD) & Documented realism assessment
(\textasciitilde60 -- 65 \%); exceeds synthetic baselines \\
No transit routing & Mode is assigned but not routed; acknowledged
limitation \\
Single peak hour & Larger tiers cover 24-hour demand patterns \\
\end{longtable}

\subsection{Construct Validity}\label{463-construct-validity}

\begin{longtable}[]{@{}ll@{}}
\toprule\noalign{}
Threat & Mitigation \\
\midrule\noalign{}
\endhead
\bottomrule\noalign{}
\endlastfoot
GEH not universally accepted & Report multiple metrics (RMSE, KS, GEH)
for triangulation \\
Reproducibility index R sensitive to μ & Handle edge cases (μ → 0)
explicitly; report CV alongside R \\
Travel time only (no queue lengths) & Focus on trip-level metrics;
acknowledge this limitation \\
\end{longtable}

\begin{center}\rule{0.5\linewidth}{0.5pt}\end{center}

\section{Reproducibility
Checklist}\label{47-reproducibility-checklist}

For any researcher to reproduce these experiments:

\begin{itemize}
\tightlist
\item[$\square$]
  Clone repository (branch \texttt{phase-14-canonical-routes} for the
  active development tip, or \texttt{main} for the thesis-tagged
  snapshot). The fully reproducible thesis-default container image is
  pinned at \texttt{ghcr.io/phanidharakula/simforge:db8d786} (see
  §3.11.5).
\item[$\square$]
  Install \texttt{uv}:
  \texttt{curl\ -LsSf\ https://astral.sh/uv/install.sh\ \textbar{}\ sh}.
\item[$\square$]
  Provision the canonical environment:
  \texttt{uv\ python\ install\ 3.13\ \&\&\ uv\ venv\ -\/-python\ 3.13\ .venv\ \&\&\ source\ .venv/bin/activate\ \&\&\ uv\ pip\ install\ -r\ requirements.lock\ \&\&\ uv\ pip\ install\ eclipse-sumo==1.26.0}
  (installs Python 3.13.13, the 35 lockfile-pinned packages including
  \texttt{path4gmns==0.10.0}, then \texttt{eclipse-sumo==1.26.0}
  separately --- the eclipse-sumo wheel is
  manylinux\_2\_28\_x86\_64-only and excluded from the cross-platform
  lockfile). Alternatively pull the pinned container; see §3.11.
\item[$\square$]
  Install Java 17+ for MATSim: \texttt{brew\ install\ openjdk@17}
  (macOS) / \texttt{apt\ install\ openjdk-17-jdk} (Linux) /
  \texttt{module\ load\ openjdk/21.0.3\_9} (Pitzer --- explicit version
  required by lmod). Then download the MATSim JAR per
  \href{../../SETUP.md}{SETUP.md}.
\item[$\square$]
  (macOS only, for DTALite OpenMP runtime)
  \texttt{brew\ install\ libomp}.
\item[$\square$]
  Fetch hash-pinned OSM PBFs: \texttt{python\ tools/download\_osm.py}
  (only required if you plan to \emph{regenerate} bundles; the committed
  \texttt{scenarios/\{chicago\_1k\_car,nyc\_10k\_car,la\_50k\_car\}/}
  networks are already built).
\item[$\square$]
  Validate the bundled scenarios:
  \texttt{python\ -m\ pipeline.validation.validate\_bundle\ scenarios/\{chicago\_1k\_car,nyc\_10k\_car,la\_50k\_car\}}.
\item[$\square$]
  Submit the canonical 55-cell benchmark on Pitzer:
  \texttt{sbatch\ cluster/jobs/benchmark\_small.sbatch}
  (\textasciitilde22-23 h walltime with Phase 12 BFS-prep cache; uses
  \texttt{-\/-mem=128G} and \texttt{-\/-cpus-per-task=16} per Phase
  12.4). On a Mac, run only the 1 K cells:
  \texttt{python\ -m\ execution.run\_benchmark\ runspecs/benchmark\_small.yaml\ -\/-scenario\ chicago\_1k\_car}.
\item[$\square$]
  (If a Pitzer worker is killed before writing its summary JSON, recover
  from disk:)
  \texttt{python\ -m\ tools.recover\_partial\_summary\ -\/-runspec\ runspecs/benchmark\_small.yaml\ -\/-scenario\ \textless{}scenario\_id\textgreater{}\ -\/-base-dir\ runs/benchmark\_small/\textless{}scenario\_id\textgreater{}}.
\item[$\square$]
  Analyse:
  \texttt{python\ -m\ evaluation.analyze\_benchmark\ runs/benchmark\_small/benchmark\_results\_benchmark\_small.json\ -\/-latex\ -\/-markdown}.
\item[$\square$]
  Audit fairness:
  \texttt{python\ -m\ evaluation.audit\_fairness\ runs/benchmark\_small}.
\item[$\square$]
  Render figures:
  \texttt{python\ -m\ evaluation.generate\_plots\ runs/benchmark\_small/benchmark\_results\_benchmark\_small.json\ -\/-output\ doc/figures}.
\item[$\square$]
  Verify: all \textasciitilde639 tests pass
  (\texttt{python\ -m\ pytest\ tests/\ -q}) with all 5 generated bundles
  present. On Apple Silicon arm64, 3 SUMO-dependent tests skip
  individually due to a known \texttt{netconvert} segfault on large
  networks; this is documented in \texttt{tests/conftest.py}.
\end{itemize}
